\scanpage{176}
\mspage{50}

Soit maintenant
\[
 M\subset\mathrm{Ob}\,\unsure{\mathcal M|M}
\]
une partie de $\mathrm{Ob}\,\mathcal M$, stable par la relation d'isomorphie. Si $M$ est $\unsure{\pi|\varpi}$-accessible \deleted{\unsure{alors}} (suivant la
condition $\mathrm b_\pi$ de $\overline{\underline{\mathrm{XIII}}}$, p. \unsure{88}), alors \deleted{\unsure{pour}}
posant
\[
 M^\pi=M\cap\mathrm{Filt}^\pi,
\]
$M^\pi$ détermine $M$ par la condition suivante
\inserted{sur $X\in\unsure{\mathrm{ob}|\mathrm{Ob}}\,\mathcal M$}
\[
\begin{aligned}
 X\in \deleted{\unsure{\mathcal M}}M\ \Longleftrightarrow\ &\text{le sous ens. }\mathrm{Ss}^M_\pi(X)\text{ de }\mathrm{Ss}_\pi(X)\\
 &\text{formé des }X_\alpha\in\mathrm{Ss}_\pi(X),\text{ tels que }X_\alpha\in M\\
 &(\Longleftrightarrow X_\alpha\in M^\pi)\text{ est cofinal dans }\mathrm{Ss}^{\deleted{\unsure{M}}}_\pi(X).
\end{aligned}\tag*{(5.1.)}
\]
Notons que lorsque $X\in\mathrm{Filt}^\pi$, ce critère est
trivialement valable, puisque alors $\mathrm{Ss}_\pi(X)$ a
$X$ comme \deleted{objet final} \inserted{plus grand élément}, et les sous-ens.
cofinaux sont ceux qui contiennent
cet objet final. Donc (5.1) est intéressant
seulement pour les $X\in\mathrm{Ob}\,\mathcal M\setminus\mathrm{Filt}^\pi$.

Notons que pour \uline{\unsure{tte|toute}} partie $M^\pi$ de
$\mathrm{Filt}^\pi$, on peut définir une partie $M$
de $\mathrm{Ob}\,\mathcal M$ par la condition (5.1.) et on
aura trivialement
\[
 M^\pi=M\cap\mathrm{Filt}^\pi.\tag*{(5.2.)}
\]
Mais sans doute, bien que $M$ satisfasse
(5.1.), il n'est pas pour \deleted{\unsure{autant}} \inserted{nécessairement}
$\pi$-accessible. En effet, pour qu'il le
soit, il \unsure{faut|doit} qu'il soit stable par $\underrightarrow{\mathrm{lim}}$

% archive sheet number: 175
